\chapter*{Hypotheses and Conventions Audit}
\phantomsection
\addcontentsline{toc}{chapter}{Hypotheses and Conventions Audit}

\section*{Riemannian and Lorentzian signatures}
Riemannian metrics are positive definite.  Lorentzian metrics use one timelike
direction; any local sign convention used in a Lorentzian chapter is stated
there explicitly.

\section*{Curvature convention}
The Riemann tensor convention is
\[
R(X,Y)Z=\nabla_X\nabla_Y Z-\nabla_Y\nabla_X Z-\nabla_{[X,Y]}Z.
\]
The Ricci tensor is the trace \( \operatorname{Ric}(Y,Z)=
\operatorname{tr}(X\mapsto R(X,Y)Z)\).  Sectional, Gaussian, scalar, and mean
curvature signs are interpreted consistently with the definitions in their
chapters.

\section*{Numerical geometry}
Discrete and numerical statements require the mesh regularity, boundary
conditions, discretization, solver tolerance, and refinement assumptions stated
in the corresponding computational chapter.  Numerical agreement is not a
substitute for the smooth theorem being approximated.
